% 结构版本 v0.2：用户已确认审阅稿（D015）；以下均为写作说明。
% 第一章 1.1–1.5 已全部定稿（V003–V010，均经用户审定）；章节级待补标记已清零（V010）。
\section{问题综述}\label{sec:problem}

\subsection{问题背景}\label{sec:background}
% 状态：V003（2026-09-25 用户对话逐句审定，1.1 定稿；段一沿用 V002 方案 A 并挂文献 hennessy2019golden / jouppi2017tpu，Crossref 核验；段二至四按版本 C 结构重写，Makespan 撤出 1.1 首现移至 1.4，D021）。素材：S01 官方题面（导语、1.1--1.3、1.5、算法优化方向）改写；不含实验数据，不提及具体方法名。

随着深度神经网络模型规模与计算复杂度持续增长，人工智能推理任务对计算吞吐、执行时延与能效提出了更高要求。神经网络处理器（Neural Processing Unit，NPU）针对神经网络计算中矩阵与向量运算密集的特点进行专用化设计，相比通用处理器能够更高效地执行推理计算\cite{hennessy2019golden,jouppi2017tpu}；然而，受制造成本、功耗与散热等因素限制，单个 AI 核心的计算能力难以持续提升。为突破单核算力瓶颈，多核 NPU 通过多个 AI 核心协同执行计算任务，成为进一步提升并行能力与能效的主流架构。

典型的多核 NPU 由多个同构的 AI 核心构成：每个核心配置矩阵计算单元（Cube）、向量计算单元（Vector）、数据搬运单元以及私有 L1/UB 缓存；所有核心共享核外主存（DDR），对 DDR 的读写共同竞争固定的总物理带宽。神经网络计算本身可表示为一张有向无环图（Directed Acyclic Graph，DAG），称为计算图：操作节点表示一次计算或数据搬运，张量节点表示计算的输入、中间结果与输出，节点间的有向边表示数据依赖。由于实际神经网络计算涉及的数据规模通常远大于单核缓存容量，执行时还需将计算拆解为可在核内执行的细粒度操作及相应的张量数据块；本题给定的计算图即以这类细粒度操作、张量及其数据依赖关系为基本描述对象。

在多核 NPU 上高效执行计算图，需要联合确定三个相互耦合的决策：将完整计算图划分为若干子图，将子图分配到各核心，并确定每个核心上的子图执行顺序。这一问题的困难源于\textbf{并行性、数据局部性与通信开销之间的固有权衡}：细化切分能够增加可独立调度的子图数量，却可能削弱数据局部性，带来跨子图边界的数据搬运与对同一输入的重复读取；将依赖关系紧密的节点聚合于同一子图，有利于中间数据的核内驻留与复用，却可能减少可独立调度的子图数量，并因子图规模增大而加重核内缓存压力，在容量不足时引发缓存换入换出。此外，各核心的核外数据搬运共同竞争有限的 DDR 带宽，局部并行度较高的调度方案并不一定取得更短的总体执行时间。

因此，多核调度本质上是一个受数据依赖、核内缓存容量、计算流水线与共享带宽共同约束的联合优化问题。进一步地，核间同步机制与共享 L2 缓存的引入会改变子图间的数据通信与缓存复用方式，对调度方法在不同硬件场景与核心数量下的适应能力提出了更高要求。如何在上述约束与硬件条件下生成合法且高效的多核切图与调度方案，是多核 NPU 高效运行所需解决的核心问题。

\subsection{多核 NPU 系统与计算图分析}
% 状态：1.2 引语段 + 1.2.1 定稿（V004，用户逐句审定，A013）；1.2.2 定稿（V007，2026-09-25 用户逐句审定＋两轮 GPT 评审仲裁，A024）；1.2.3 定稿（V006，用户逐句审定，A021）。
为明确后续建模涉及的硬件约束、计算对象与调度方案表示，本节首先介绍多核 NPU 的硬件架构，随后说明计算图的结构与节点表示，并给出子图划分与任务调度的组织方式。

\subsubsection{多核 NPU 硬件架构}
% 状态：V004 定稿（2026-09-25 用户逐句审定）。素材：题面 1.1/1.2/1.5、附录 A、B.4、D.5；无实验数据、无方法名；不挂文献（"商用一致"断言过强已撤，Ascend 回 D024 候选池）。
题设多核 NPU 的核心执行资源包括四类执行管道：PIPE\_M 与 PIPE\_V 分别承载 Cube 矩阵计算与 Vector 向量计算；数据搬运管道 PIPE\_MTE2 负责 DDR$\rightarrow$L1/UB 与 UB$\rightarrow$L1 方向的搬运，PIPE\_MTE3 负责 L1/UB$\rightarrow$DDR 与 L1$\rightarrow$UB 方向的搬运。在官方评估模型中，同一管道上的操作串行执行，不同管道上满足数据依赖与资源约束的操作可并行执行，因此 Cube 与 Vector 可同时工作。

存储层级由核内私有缓存与核外共享主存构成：每个核心拥有私有 L1 缓存（524288 bytes，合 512 KiB）与 UB（131072 bytes，合 128 KiB），分别主要服务于 Cube 与 Vector 计算；所有核心共享的 DDR 提供固定的总物理带宽（60 bytes/cycle），由 DDR 服务的核外数据搬运共同竞争该带宽，而 L1 与 UB 之间的核内搬运只占用对应管道、不消耗 DDR 带宽。DDR 访问不按固定周期数结算，完成时间由搬运数据量与评估时刻的可用带宽共同决定。此外，问题三在上述结构基础上引入全核共享的只读 L2 缓存（容量 1 MB、读带宽 250 bytes/cycle，与 DDR 带宽相互独立），用于复用多核共享输入，其缓存机制与调度含义在第六章展开。多核 NPU 的固定硬件配置汇总于表~\ref{tab:hw-config}。

\begin{table}[htbp]
  \centering
  \caption{多核 NPU 固定硬件配置}
  \label{tab:hw-config}
  \widetable{\begin{tabular}{llp{5.9cm}}
    \toprule
    部件 & 参数 & 数值 \\
    \midrule
    L1 缓存（每核私有） & 容量 & 524288 bytes（512 KiB） \\
    UB（每核私有） & 容量 & 131072 bytes（128 KiB） \\
    DDR（全核共享） & 总物理带宽 & 60 bytes/cycle \\
    L2 只读缓存（全核共享，仅问题三） & 容量 / 读带宽 & 1048576 bytes（题面记为 1 MB）/ 250 bytes/cycle \\
    \bottomrule
  \end{tabular}}
  \begin{minipage}{0.9\linewidth}
    \small 注：1 KiB = 1024 bytes；取值来自题面附录 A 及 B.4 固定配置。
  \end{minipage}
\end{table}

% 图 F05（1.2.1 配图）：用户以 PowerPoint 手工绘制的终图 figs/F05_hw_abstraction.png（pptx 源文件待归档 figs/）；matplotlib 参考版与脚本见 figs/F05_hw_abstraction_scriptref.* 与 scripts/fig_F05_hw_abstraction.py（来源登记见 05/figure_registry、09、08/D034）。
\begin{figure}[htbp]
  \centering
  \includegraphics[width=0.92\textwidth]{F05_hw_abstraction.png}
  \caption{多核 NPU 硬件抽象与存储层级示意图}
  \label{fig:hw-arch}
\end{figure}

\subsubsection{计算图结构与节点表示}
% 状态：V007 定稿（2026-09-25 用户逐句审定＋两轮 GPT 评审仲裁，A024）。素材：题面 1.2、附录 B.3/B.7；无实验数据、无方法名、不挂文献。X01：正文按题设文字口径并加"按照题设给定的计算图格式"限定（官方 stub 代码 _build_op_adjacency 允许 Op→Op 直连边，T014 核验后如证实全量输入满足可去限定）；s_t 的跨边界搬运语义归 1.2.3 评估器句。1.2.3 字段名同轮统一 \texttt（纯排版，D031 不占号）。
% 图位（F20 输入计算图最小示例；草图 v3.1 已交待审核，tex 插图随 F06/F20 插图轮进行，登记见 registry/05）：以题面附录 B.7 最小例（DDR Tensor→COPY_IN→UB Tensor→ADD→UB Tensor→COPY_OUT→DDR Tensor）展示两类节点、边方向与 pos 属性；图管关系、表管字段，节点只标代表性属性（Tensor: pos/size，Op: op/pipe/cycles），不复刻题面图 1 矩阵乘法例。
如 1.1 节所述，本题给定的计算图以细粒度操作、张量及其数据依赖关系为基本描述对象，整体构成有向无环图。形式上，将计算图记为 $G=(V,E)$：张量节点集合 $T$ 与操作节点集合 $O$ 不相交，共同构成节点集合 $V=T\cup O$，每个节点具有在同一计算图内全局唯一的整数 \texttt{id}。按照题设给定的计算图格式，有向边集合 $E$ 只允许张量节点与操作节点相连：对 $t\in T$、$o\in O$，边 $t\rightarrow o$ 表示张量 $t$ 是操作 $o$ 的输入，边 $o\rightarrow t$ 表示 $t$ 是 $o$ 的输出。由此，计算图具有张量与操作相连的二部结构（Tensor--Op 二部图），操作之间的数据依赖经由其间的张量传递；边本身不设置独立的数据量字段，相关数据规模由关联张量的 \texttt{size} 属性给出。

张量节点（Tensor）表示一块数据，记作 $t\in T$，可作为计算图的输入、输出或中间结果，包含位置与数据大小两类属性。位置属性 \texttt{pos} 表示输入图规定的逻辑存储位置，可取 DDR、L1 或 UB：计算图的输入、输出张量位于 DDR，中间张量位于核内缓存（L1 或 UB）。数据大小 $s_t$（字段 \texttt{size}）为张量占用的字节数，取非负整数，同时用于核内容量检查与数据搬运量计算。

操作节点（Op）表示一次计算或数据搬运，记作 $o\in O$，包含操作类型、执行管道与周期字段三类属性。操作类型 \texttt{op} 标识具体操作：COPY\_IN 将数据从 DDR 搬入核内缓存，COPY\_OUT 将数据从核内缓存搬回 DDR，ADD、MUL 等为在指定计算单元上执行的核内计算操作。执行管道 \texttt{pipe} 指定操作占用的硬件流水线，取 1.2.1 节所述四类执行管道之一：核内计算操作在 PIPE\_M 或 PIPE\_V 上执行，搬运操作占用相应的数据搬运管道。对核内计算操作，字段 \texttt{cycles} 给出其执行周期数 $c_o$，由计算图直接给定；对于访问 DDR 的搬运操作，其最终耗时不由 \texttt{cycles} 直接确定，而由搬运数据量与评估时刻的可用带宽共同决定（见 1.2.1）。

输入计算图采用 JSON 格式，顶层包含 \texttt{tensors}、\texttt{ops} 与 \texttt{edges} 三个数组，分别记录张量节点、操作节点与有向边，各字段含义汇总于表~\ref{tab:input-format}。题设对输入计算图进一步保证：每个 DDR 输入只搬入一次、每个最终结果只搬出一次；中间计算不与 DDR 交换数据；单个计算操作的输入、输出总量不超过相应核内缓存容量。

\begin{table}[htbp]
  \centering
  \caption{输入计算图 JSON 字段说明}
  \label{tab:input-format}
  \widetable{\begin{tabular}{llp{5.9cm}}
    \toprule
    数组 & 字段 & 含义与取值 \\
    \midrule
    \texttt{tensors} & \texttt{id} & 张量节点 id，与其他节点共享全局唯一编号 \\
     & \texttt{pos} & 逻辑存储位置：DDR / L1 / UB \\
     & \texttt{size} & 张量大小 $s_t$（bytes，非负整数） \\
    \texttt{ops} & \texttt{id} & 操作节点 id，与其他节点共享全局唯一编号 \\
     & \texttt{op} & 操作类型：COPY\_IN / COPY\_OUT / ADD / MUL 等 \\
     & \texttt{pipe} & 执行管道：PIPE\_M / PIPE\_V / PIPE\_MTE2 / PIPE\_MTE3 \\
     & \texttt{cycles} & 周期字段；核内计算操作据此确定执行周期 $c_o$ \\
    \texttt{edges} & \texttt{source} / \texttt{target} & 边两端节点 id；题设规定为 Tensor\,$\rightarrow$\,Op 或 Op\,$\rightarrow$\,Tensor \\
    \bottomrule
  \end{tabular}}
  \begin{minipage}{0.9\linewidth}
    \small 注：字段与约束来自题面附录 B.3；数值字段均为非负整数，所有节点 \texttt{id} 在同一计算图内全局唯一。
  \end{minipage}
\end{table}

\subsubsection{子图划分与任务调度描述}
% 状态：V006 定稿（2026-09-25 用户逐句审定，A021）。素材：题面 1.2/1.3/1.5、附录 B.5/B.7/D.1/D.2；无实验数据、无方法名、不挂文献。延后项（V009 轮更新去向）：延迟常量（100/1000/500 cycles）不入 1.5，留第四至六章对应模型与实验设置；$C_{\mathrm{inner}}/C_{\mathrm{cross}}$ 仅在后文确有使用时引入；Makespan 等指标名不提前出现（D021 首现 1.4）；场景 B 驻留代价与容量影响按拆层落实——问题定义层在 1.5.2（V009），驻留机理与权衡分析留第五章（跨节去重）。COPY\_IN/COPY\_OUT 为 1.2.2 待定义节点类型，1.2.2 落稿时衔接。
在计算图表示的基础上，多核调度按两个相互衔接的过程展开：切图通过确定核内操作节点的子图归属，将完整计算图划分为若干子图，子图数量不作限制；子图调度将各子图分配到具体核心，并确定每个核心上的子图执行顺序。划分得到的每个子图以子图 id（Subgraph id，下文简称 sgid）标识。切图方案以映射表 \texttt{node\_to\_subgraph} 表示：核内操作节点（即除 COPY\_IN 与 COPY\_OUT 两类数据搬运节点外的操作节点）各自映射到一个非负整数 sgid，且映射必须覆盖全部核内操作节点。跨子图的依赖边在子图层面构成子图依赖图；子图依赖图必须为有向无环图，即至少存在一个合法的拓扑序，否则切图方案无效。

子图在核心上的执行以任务（Task）为单位组织：Task 是 CPU 主控向单个核心一次下发并统一调度的最小执行单元，一个 Task 可包含一个或多个子图，对应关系由硬件场景决定。场景 A 无核间同步机制，每个子图独立构成一个 Task；Task 之间不保留核内缓存状态，跨子图数据都必须经 DDR 中转。场景 B 引入核间同步机制：在官方评估规则下，同一核心上的全部子图合并为一个 Task，前序子图的输出可驻留在 L1/UB 中供后续子图直接使用，Task 内的子图边界不再插入数据搬运；跨核依赖则由源核心将数据写入 DDR 并发送同步信号，目标核心收到信号后读取。问题三沿用场景 B 的 Task 组织方式，仅增加全核共享的只读 L2 缓存（见 1.2.1）。

子图调度方案以 \texttt{core\_schedules} 表示：由 $N$ 个 sgid 列表组成（$N$ 为核心数量），外层下标对应核心编号，内层列表按执行顺序依次给出该核心分得的 sgid，允许为空，即允许个别核心不分得任何子图；全部列表合并后，每个 sgid 必须且只能出现一次（即每个子图恰好被调度一次），同一核心内的执行顺序还不得违反子图间依赖，否则子图调度方案无效。方案本身只给出子图层面的归属、分核与顺序；官方评估程序根据原始计算图、切图方案、子图调度方案及对应场景规则，重建跨 Task 边界的数据搬运，再确定各 Task 内部的核内操作顺序，并在共享带宽约束下模拟多核并行执行，最终输出各项评估指标。

\subsection{输入图解析与调度特征提取}
% 状态：V010 定稿（2026-09-26 用户逐段审定＋三轮 GPT 评审仲裁，AI-09/A046）。素材：官方评估器与 stub 代码核验（校验/调度两条建图视图、合法性两层、P3 缓存键=COPY_IN 输出张量 logical\_tid 跨核共享）、v5 两段式实现核料（实际特征清单＋"未找到"清单：无独立输入校验/无静态容量约束/接受准则只看 Makespan）、2025A 国一 1.2/1.3 写法对照（解析入一章、特征公式后移）；无实验数据、无方法名、不挂文献。要点：三层信息框架（直接读取／静态推导／评估获得或统计）贯穿全节；COPY 收缩细节与"与官方评估程序一致"整称断言撤出（求解侧辅助索引＋官方为准）；L2 复用对象不限定 DDR 图输入（同一张量跨核重复 COPY\_IN 共享缓存键）；负载项章节指向用"后续"（三问共用）。1.3 无图表（D045 复核）。延后项：分管道计算量等特征的正式符号留第四至六章首用处定义（D039 符号三处分工）；术语候选（多消费者张量/分管道计算量）待第四章与 2.1 使用时对齐入 06（Needs-integration）。
1.2 节说明了计算图与调度方案的表示。在此基础上，本文求解流程对输入计算图完成解析与信息提取：先由输入文件构建图结构与依赖关系（1.3.1 节），再从中得到支撑切图与调度构造、评价的各类信息（1.3.2、1.3.3 节）。按照获取方式，这些信息分为三层——直接读取的输入字段属性、由图结构与候选方案静态推导的特征、只能由评估过程获得或基于评估结果统计的指标。

\subsubsection{图结构解析}
求解流程读入 1.2.2 节所述 JSON 格式的输入计算图，为张量节点与操作节点分别建立属性表，并由边表建立张量与操作之间的生产、消费关联：指向操作节点的边将其记为该操作的输入张量，由操作节点引出的边将其记为该操作的输出张量。在张量—操作关联的基础上，为参与切图的核内操作节点（见 1.2.3 节）建立前驱与后继依赖索引，作为后续结构分析与候选方案构造的处理基础。该索引是求解侧的辅助表示，不改变输入计算图本身；方案的合法性判定与正式执行过程始终以官方评估程序为准。

依托该依赖索引可获得核内操作层面满足数据依赖约束的拓扑次序；输入计算图本身的无环性由官方评估入口统一校验。输入图与候选方案的合法性均由对应官方评估入口统一检查：基础校验涵盖输入字段类型与取值、节点编号全局唯一、边端点合法以及 1.2.3 节所列的方案条件，场景相关的执行可行性由对应问题的评估流程继续检查；求解流程不另设与官方并行的判定口径。候选切图方案产生后，将核内操作间的依赖按 \texttt{node\_to\_subgraph} 映射聚合：两端分属不同子图的依赖收缩为对应子图间的依赖边，由此构建子图依赖图，并以拓扑排序检验其无环性。

\subsubsection{计算与通信特征提取}
第一层信息由输入字段直接读取。1.2.2 节表~\ref{tab:input-format} 所列属性在调度中各司其职：字段 \texttt{cycles} 给出的执行周期 $c_o$ 度量核内计算操作的开销，是各类计算量统计的基本单元；字段 \texttt{size} 给出的张量大小 $s_t$ 以字节计量数据规模，是搬运量与存储占用的基本依据；执行管道 \texttt{pipe} 区分操作占用的资源维度，使负载可按管道分别核算；位置属性 \texttt{pos} 区分核内与核外数据；操作类型区分参与划分的核内操作与承担边界搬运的两类 COPY 操作。字段含义不再复述，以下只说明在其之上构造的调度特征。

第二层特征由图结构与候选分组静态推导，不经评估即可算得，对任意给定的操作分组均可计算，候选切图方案下即取各子图。分管道计算量将组内操作按执行管道分别累计执行周期（单位 cycles）；由于 Cube 与 Vector 两类计算管道可并行执行（1.2.1 节），计算负载须分管道核算而非合并为单一总量，该特征构成后续各核心负载估计的基础。跨组生产—消费张量字节把生产方与消费方分属不同分组的张量，按其大小累加到分组对上（单位 bytes），刻画候选划分形成的潜在边界数据规模；当核心分配形成跨 Task 或跨核边界时，该量可用于估计 DDR 通信需求，场景 B 下同核的生产—消费对可经核内驻留复用避免搬运（1.2.3 节），实际额外搬运量仍以官方评估结果为准。多消费者张量记录被两个及以上分组消费的张量的大小与消费方集合，为跨组通信与潜在重复读取分析提供结构信息，其与共享 L2 缓存复用的关联见 1.3.3 节。

\subsubsection{缓存相关特征提取}
与核内存储相关的信息来自张量属性与依赖结构。位置属性把张量分为核内驻留（L1 或 UB）与核外（DDR）两类，张量大小给出占用字节数；由生产—消费关系可确定各张量的产生端、消费端与复用关系，为判断潜在驻留范围与缓存压力提供结构信息。题设已保证单个操作的输入输出总量不超过相应核内缓存容量（1.2.2 节）。对子图或 Task 整体的实际驻留峰值与缓存换入换出（spill），本文不作静态近似：二者由官方评估程序的核内调度在评估中确定。因此，本文不把分组内张量大小之和当作容量约束使用；容量相关影响经由评估结果体现——spill 产生的搬运计入总额外数据搬运量，并通过对总体执行时间的作用进入候选方案的评估比较。

与共享 L2 缓存相关的静态访问特征是同一张量跨核重复读取的识别：当核心分配使多个核心对同一张量各自发起 COPY\_IN 访问时，该张量成为潜在的复用对象，其字节量与消费方分布（即 1.3.2 节的多消费者张量信息）刻画可复用空间；其与切图、分配及执行顺序的关联见 1.5.3 节，缓存机制与利用分析见第六章，Cache 命中率与命中字节由评估过程反馈获得。综上，本文调度的信息基础分为三层：直接读取的输入字段属性（1.2.2 节）、由图结构与候选方案静态推导的特征（1.3.2 节与本节），以及评估获得或基于评估结果进一步统计的性能指标——总体任务执行时间、总额外数据搬运量、Cache 命中率及加速比（1.4 节）。前两层支撑候选方案的构造与静态代价刻画，第三层用于方案评价与结果统计；其中总体任务执行时间作为候选筛选与方案改进的直接评价依据，其余指标用于辅助分析和结果报告。各类信息在第四至六章中的具体使用方式随相应模型与算法展开。

\subsection{名词解释}
% 状态：V008 定稿（2026-09-25 用户逐句审定＋四轮 GPT 评审仲裁，A029）。素材：题面 1.4 评估指标整节、问题一加速比定义与脚注 [^2][^3]、1.2 各定稿节；遵循 D036 词条自足原则（定义可复述、机制不重复）。符号首现：S_{k,i}、T_{1,i}、T_{k,i}、\bar S_k、m，3.2 符号表（TBL-003）届时收录；单核基准构造（启发式整图核内调度固定结果）与 P3 相对加速比定义留实验章/第六章；Makespan 与总额外数据搬运量首现本节（D021）。
本节对后文频繁使用的主要术语与评价指标作统一说明；相关硬件结构、计算图表示与调度方案的机制细节见 1.2 节，各评价指标在第四至六章的实验中按本节口径统计与报告。

\begin{itemize}
\item \textbf{NPU（神经网络处理器）}：面向神经网络计算中矩阵与向量运算密集的特点进行专用化设计的处理器；本题平台由 $N$ 个同构 AI 核心构成多核 NPU，正文中的“核心”均指 AI 核心。
\item \textbf{计算图（DAG）}：描述神经网络计算的有向无环图，操作节点与张量节点经有向边连接，边方向表达数据依赖。切图与子图调度均以计算图为直接对象。
\item \textbf{Tensor（张量节点）}：计算图中承载数据的节点，可表示输入、输出或中间结果。张量的数据大小是核内容量检查与数据搬运量计算的基本依据。
\item \textbf{Op（操作节点）}：计算图中表示一次计算或数据搬运的节点。核内计算操作（如 ADD、MUL）在相应计算单元上执行，COPY\_IN 与 COPY\_OUT 两类搬运操作负责核内缓存与 DDR 之间的数据交换。
\item \textbf{Pipe（执行管道）}：操作执行时占用的硬件流水线，本题包括 PIPE\_M、PIPE\_V、PIPE\_MTE2 与 PIPE\_MTE3 四类。在官方评估模型中，同一管道上的操作串行执行，不同管道上的操作在满足依赖与资源约束时可并行。
\item \textbf{Subgraph（子图）}：切图得到的核内操作节点分组，是核心分配与执行顺序的基本单位，子图数量不作限制。COPY\_IN 与 COPY\_OUT 两类搬运节点不参与子图分配。
\item \textbf{sgid（子图 id）}：子图的标识，取非负整数。切图方案以核内操作节点到 sgid 的映射表达，子图调度方案以各核心上的 sgid 有序列表表达。
\item \textbf{Task（任务）}：CPU 主控向单个核心一次下发并统一调度的最小执行单元，可包含一个或多个子图：场景 A 中每个子图独立构成一个 Task，场景 B 中同一核心上的全部子图合并为一个 Task。Task 的组织方式影响跨子图数据能否在核内驻留复用以及是否需要经 DDR 搬运。
\item \textbf{缓存换入换出（spill）}：官方评估程序在核内调度阶段因 L1/UB 容量不足自动插入的搬运：换出将驻留数据写回 DDR，换入在其再次使用前搬回。二者消耗 DDR 带宽并计入额外数据搬运量。
\end{itemize}

在上述术语的基础上，本文采用以下指标评价多核切图与调度方案的性能。

\begin{itemize}
\item \textbf{总体任务执行时间（Makespan）}：最重要的评估指标，指官方评估程序完成多核模拟后，所有操作结束时刻的最大值，单位为时钟周期。
\item \textbf{总额外数据搬运量}：次要指标，指相对于原始计算图已有 DDR 搬运、由切图与调度额外产生的 DDR 搬运字节数；新增搬运主要来自跨 Task 边界数据、对同一输入的重复读取与缓存换入换出。
\item \textbf{加速比}：对测试用例 $i$ 的 $k$ 核配置，定义
  \[ S_{k,i}=\frac{T_{1,i}}{T_{k,i}} \]
  其中 $T_{1,i}$ 与 $T_{k,i}$ 分别为该用例的单核基准 Makespan 与 $k$ 核多核评估所得 Makespan，单核加速比定义为 1。对包含 $m$ 个测试用例的测试集，平均加速比定义为
  \[ \bar{S}_k=\frac{1}{m}\sum_{i=1}^{m} S_{k,i} \]
  即取各测试用例加速比的算术平均，而非单核总时间与多核总时间之比。问题三另使用相对无 L2 基线的加速比，见第六章。
\item \textbf{Cache 命中率（仅问题三）}：次要指标，指可由只读 Cache 服务的 COPY\_IN 访问中，命中字节数占总访问字节数的比例，用于观察只读 Cache 对多核共享输入访问的覆盖程度。
\end{itemize}

\subsection{拟解决的问题}\label{sec:problems}
% 状态：V009 定稿（2026-09-26 用户对话审定＋四轮 GPT 评审仲裁，A036）。写法对照 2025A 国一 1.5：取其职责单一纪律（只答"解什么、题面要什么结果"），按冻结大纲改为总起段＋TBL-001＋三个单段小节；"为什么难"归第二章、驻留机理归第五章、结果展示集中表内（07 窗口表 W-A"每问三段式"提示被本定稿取代，用户确认）。素材：题面 1.3、1.5 硬件场景、问题一～三题述与脚注 1、附录 B.6/D.2；无实验数据、无方法名、不挂文献。
在 1.2 节的计算图与调度方案表示、1.4 节的评价指标基础上，本题要求提出高效且效果良好的多核切图与调度方法：对给定计算图与核心数量，联合确定各核内操作节点的子图归属、子图的核心分配以及各核心上的子图执行顺序，在合理时间内生成合法的切图方案与子图调度方案，以\textbf{缩短总体任务执行时间}为主要目标并控制额外数据搬运。三个问题在方案层面直接输出的决策对象相同，差异主要来自硬件场景及对应的评估机制：问题一考虑无核间同步机制的场景 A；问题二引入核间同步机制，即场景 B，Task 组织方式与核内数据复用方式随之改变；问题三在场景 B 基础上引入全核共享的只读 L2 缓存。三问均要求算法适应核数变化、稳定而高效地生成高质量结果，其中问题一、二明确要求覆盖 2～5 核（“高效”指避免暴力迭代搜索，以单个用例在 5～10 分钟内给出结果为佳）。实验均在题目给定的全部测试用例和固定配置参数下完成；三问的硬件场景、核心任务与题面结果要求汇总于表~\ref{tab:problem-reqs}。

% 表 TBL-001（三问场景与交付对照表）：V009 首次成表，四列定版（问题/硬件场景/核心任务/题面结果要求，05 与 registry 已同轮同步列定义，D029）；交付说明只在本表出现，三个小节不再复述。与 2.5 TBL-002 分工（D042 编号顺延）：本表管题面任务与结果要求，机制差异与决策边界归 TBL-002（05 §"TBL-002 不并入 TBL-001"）。
\begin{table}[htbp]
  \centering
  \caption{三个问题的硬件场景、核心任务与题面结果要求}
  \label{tab:problem-reqs}
  \widetable{\begin{tabular}{lp{3.4cm}p{3.4cm}p{3.9cm}}
    \toprule
    问题 & 硬件场景 & 核心任务 & 题面结果要求 \\
    \midrule
    问题一 & 场景 A：无核间同步，每个子图独立构成一个 Task，跨子图数据均经 DDR 中转 & 建立场景 A 下的多核切图与调度模型，以缩短总体任务执行时间为主、兼顾总额外数据搬运量 & 正文：1～5 核平均加速比折线图；附录：逐用例 Makespan 与总额外数据搬运量；附件：可复现 Python 程序脚本 \\
    问题二 & 场景 B：引入核间同步，同核全部子图合并为一个 Task，同 Task 内直接复用核内驻留数据，跨核经 DDR 中转与同步信号 & 在问题一基础上，在严格满足 L1/UB 容量约束的前提下进一步缩短总体任务执行时间并减少额外数据搬运 & 与问题一相同（场景 B 下） \\
    问题三 & 在场景 B 基础上增加全核共享只读 L2 缓存（复用多核共享输入），Task 组织与问题二相同 & 在题目给定配置下，研究引入共享只读 L2 对多核切图与调度的影响，并建立相应模型 & 正文：无 L2 与只读 Cache 两种配置的 1～5 核对比曲线及相对无 L2 基线的加速比；附录：两配置逐用例 Makespan、总额外数据搬运量与 Cache 命中率 \\
    \bottomrule
  \end{tabular}}
  \begin{minipage}{0.9\linewidth}
    \small 注：三问在方案层面直接输出的决策对象相同，均为切图方案与子图调度方案（见 1.2.3 节）；本表依据题面问题一至三的要求整理。
  \end{minipage}
\end{table}

\subsubsection{问题一：场景 A 下的多核切图与调度}
% 状态：V009 定稿（同上）。要点：场景 A 事实层——每子图一 Task、跨子图必经 DDR 且不因是否同核而改变（附录 D.2）；不写"任意两子图单位通信成本相同"，避免与 Task 激活等待 100/1000 cycles 的不对称产生表面冲突；约束并列删除（合法性由"1.2.3 节全部合法性条件"承担，流水线/带宽属评估机制）；效率与实验要求只在总起段出现一次；交付只在表内。
如 1.2.3 节所述，场景 A 无核间同步机制：每个子图独立构成一个 Task，Task 间不保留核内缓存状态，跨子图数据均需经 DDR 中转，不因两个子图是否分配在同一核心而改变。问题一要求在此场景下建立多核切图与调度模型并设计求解算法：给定输入计算图与核心数量 $N$，为每个核内操作节点确定子图归属，并确定各子图的核心分配与各核心上的执行顺序，输出满足 1.2.3 节全部合法性条件的切图方案与子图调度方案，以缩短总体任务执行时间为主要目标，兼顾总额外数据搬运量。

\subsubsection{问题二：场景 B 下的多核切图与调度}
% 状态：V009 定稿（同上）。要点：方案表示与基本合法性条件沿用问题一；场景 B 三事实（合并 Task 带"在官方评估规则下"限定／同 Task 内复用／跨核 DDR＋同步信号）；驻留代价拆层——本节只保留问题定义层"复用避免边界搬运但驻留占用有限 L1/UB 空间"（引出严格容量约束），机理链（压缩中间子图空间、降低执行效率）与权衡分析留第五章，第二章是否先概括待其起草定（1.2.3 注记"留 1.5 问题二/第五章"在拆层下仍为真）；交付只在表内。
问题二将硬件场景由场景 A 换为场景 B，方案表示与基本合法性条件沿用问题一。场景 B 引入核间同步机制：如 1.2.3 节所述，同一核心上的全部子图在官方评估规则下合并为一个 Task，同 Task 内的子图可直接复用核内驻留数据，跨核依赖仍经 DDR 中转与同步信号传递。同核复用能够避免同 Task 内子图边界的数据搬运，但驻留数据持续占用有限的 L1/UB 空间；问题二因而要求在问题一的基础上建立适用于场景 B 的切图与调度模型，在\textbf{严格满足 L1/UB 容量约束}的前提下，进一步缩短总体任务执行时间并减少额外数据搬运。

\subsubsection{问题三：共享 L2 Cache 下的多核切图与调度}
% 状态：V009 定稿（同上）。要点：因果方向取题面原语——研究共享只读 L2 对多核切图与调度的影响（不加"性能"）；调度→访问模式→Cache 复用效果→时间为机理句而非任务定义；"题目给定的容量与带宽配置"限定（不解读为参数扫描）；两配置对照是任务组织方式留在正文，字段清单归表；P3 结尾不宣布额外搬运为优化目标（题面只要求报告）；补"设计求解算法"与 P1/P2 对仗。
问题三在场景 B 的基础上引入全核共享的只读 L2 缓存（见 1.2.1），用于复用多核共享输入，Task 组织方式与问题二相同，决策对象与合法性条件不变。该问题要求在题目给定的只读 L2 容量与带宽配置下，\textbf{研究共享只读 L2 对多核切图与调度的影响}，并建立相应模型、设计求解算法；切图、核心分配与执行顺序会改变共享输入的访问模式，进而影响 Cache 复用效果及总体执行时间。其任务以无 L2 与只读 Cache 两种配置的对照形式组织，用以比较共享只读 L2 对总体任务执行时间、总额外数据搬运量及 Cache 复用效果的影响。
